\documentclass[12pt,a4paper]{article}
\usepackage[T1]{fontenc}
\usepackage{amsmath,amssymb,mathrsfs,tikz,times,pifont}
\usepackage{enumitem}
\usepackage{multicol}
\usepackage{lmodern}
\usetikzlibrary{trees}
\newcommand\circitem[1]{%
\tikz[baseline=(char.base)]{
\node[circle,draw=gray, fill=red!55,
minimum size=1.2em,inner sep=0] (char) {#1};}}
\newcommand\boxitem[1]{%
\tikz[baseline=(char.base)]{
\node[fill=cyan,
minimum size=1.2em,inner sep=0] (char) {#1};}}
\setlist[enumerate,1]{label=\protect\circitem{\arabic*}}
\setlist[enumerate,2]{label=\protect\boxitem{\alph*}}
%%%::::::by chnini ameur :::::::%%%
\everymath{\displaystyle}
\usepackage[left=1cm,right=1cm,top=1cm,bottom=1.7cm]{geometry}
\usepackage[colorlinks=true, linkcolor=blue, urlcolor=blue, citecolor=blue]{hyperref}
\usepackage{array,multirow}
\usepackage[most]{tcolorbox}
\usepackage{varwidth}
\usepackage{float} %pour utiliser l'option [H] qui force l'image à apparaître exactement à l'endroit où elle est placée dans le code.
\tcbuselibrary{skins,hooks}
\usetikzlibrary{patterns}
%%%::::::by chnini ameur :::::::%%%
\newtcolorbox{exa}[2][]{enhanced,breakable,before skip=2mm,after skip=5mm,
colback=yellow!20!white,colframe=black!20!blue,boxrule=0.5mm,
attach boxed title to top left ={xshift=0.6cm,yshift*=1mm-\tcboxedtitleheight},
fonttitle=\bfseries,
title={#2},#1,
% varwidth boxed title*=-3cm,
boxed title style={frame code={
\path[fill=tcbcolback!30!black]
([yshift=-1mm,xshift=-1mm]frame.north west)
arc[start angle=0,end angle=180,radius=1mm]
([yshift=-1mm,xshift=1mm]frame.north east)
arc[start angle=180,end angle=0,radius=1mm];
\path[left color=tcbcolback!60!black,right color = tcbcolback!60!black,
middle color = tcbcolback!80!black]
([xshift=-2mm]frame.north west) -- ([xshift=2mm]frame.north east)
[rounded corners=1mm]-- ([xshift=1mm,yshift=-1mm]frame.north east)
-- (frame.south east) -- (frame.south west)
-- ([xshift=-1mm,yshift=-1mm]frame.north west)
[sharp corners]-- cycle;
},interior engine=empty,
},interior style={top color=yellow!5}}
%%%%%%%%%%%%%%%%%%%%%%%
\usepackage{fancyhdr}
\usepackage{eso-pic}         % Pour ajouter des éléments en arrière-plan
% Commande pour ajouter du texte en arrière-plan
\usepackage{tkz-tab}
\AddToShipoutPicture{
    \AtTextCenter{%
        \makebox[0pt]{\rotatebox{80}{\textcolor[gray]{0.7}{\fontsize{5cm}{5cm}\selectfont PGB}}}
    }
}
\usepackage{lastpage}
\fancyhf{}
\pagestyle{fancy}
\renewcommand{\footrulewidth}{1pt}
\renewcommand{\headrulewidth}{0pt}
\renewcommand{\footruleskip}{10pt}
\fancyfoot[R]{
\color{blue}\ding{45}\ \textbf{2026}
}
\fancyfoot[L]{
\color{blue}\ding{45}\ \textbf{Prof:M. BA}
}
\cfoot{\bf
\thepage /
\pageref{LastPage}}
% Création du compteur pour les exercices
\newcounter{exercice}
\renewcommand{\theexercice}{\arabic{exercice}}  % Définit l'affichage du compteur en chiffres arabes

% Définir la commande \exo
\newcommand{\exo}{\refstepcounter{exercice}\textbf{Exercice \theexercice} }
\begin{document}
\renewcommand{\arraystretch}{1.5}
\renewcommand{\arrayrulewidth}{1.2pt}
\begin{tikzpicture}[overlay,remember picture]
    \node[draw=blue,line width=1.2pt,fill=purple,text=blue,inner sep=3mm,rounded corners,pattern=dots]at ([yshift=-2.5cm]current page.north) {\begingroup\setlength{\fboxsep}{0pt}\colorbox{white}{\begin{tabular}{|*1{>{\centering \arraybackslash}p{0.28\textwidth}} |*2{>{\centering \arraybackslash}p{0.2\textwidth}|} *1{>{\centering \arraybackslash}p{0.19\textwidth}|} }
                \hline
                \multicolumn{3}{|c|}{$\diamond$$\diamond$$\diamond$\ \textbf{Lycée de Dindéfélo}\ $\diamond$$\diamond$$\diamond$ } & \textbf{A.S. : 2026/2027}                                              \\ \hline
                \textbf{Matière: Mathématiques}                                                                                    & \textbf{Niveau : 2}\textbf{ndL} & \textbf{Date: 17/11/2026} & \textbf{} \\ \hline
                \multicolumn{4}{|c|}{\parbox[c]{10cm}{\begin{center}
                                                                  \textbf{{\Large\sffamily Td Calcul Dans $\mathbb{R}$}}
                                                              \end{center}}}                                                                                                        \\ \hline
            \end{tabular}}\endgroup};
\end{tikzpicture}

\vspace{3cm}
\fbox{\textbf{\exo}}

Ecrire les réels suivants sous formes de fractions irréductibles :
\begin{multicols}{2}
\setlength{\columnseprule}{0.1mm}
\begin{enumerate}
\item $\textbf{A}=\frac{\frac{3}{2}-\frac{5}{3}}{-\frac{7}{2}+\frac{1}{6}} \times \frac{1}{-3+\frac{2}{5}}$
\item $\textbf{B}=\frac{2}{3}-\frac{4}{5}+\frac{3}{7}$
\item $\textbf{C}=\left(3+\frac{4}{5} \right) \left(\frac{1}{5}-4 \right)$
\item $\textbf{D}=\left(\frac{2}{3}-1 \right)\times \frac{2+\frac{1}{3}}{\frac{3}{5}+\frac{1}{5}}$
\item $\textbf{E}=\frac{3+\frac{1}{3}}{\frac{2}{3}-2} \times \frac{2}{\frac{1}{3}+3}$
\item $\textbf{F}=\frac{\frac{1}{4}+3}{\frac{3}{5}} \times \frac{7}{3}$
\item $\textbf{G}= \frac{5}{2}-\frac{1}{\frac{2}{3}}$
\item $\textbf{H}= \frac{\frac{2}{3}-\frac{2}{3}}{\frac{5}{4}} - \frac{7}{15}$
\item $\textbf{I}= \frac{\frac{1}{2}+\frac{2}{3}-1}{\frac{3}{4}}$
\item $\textbf{J}= \left(5-\frac{3}{2} \right)  \left(\frac{3}{4}+\frac{-5}{4} \right)$
\item $\textbf{K}=\left(\frac{5}{3}+1 \right) \div \left(\frac{4}{3}-\frac{1}{4} \right)  $
\item $\textbf{L}=\left(3-\frac{2}{3} \right) + \left(-\frac{3}{2}+3 \right)   $
\end{enumerate}
\end{multicols}

\fbox{\textbf{\exo}}

\begin{multicols}{2}
\setlength{\columnseprule}{0.1mm}
\begin{enumerate}
\item Développer et réduire les expressions données :
\begin{enumerate}
\item $A = (-5x - 4)(2x^2 + x - 3) + (2x - 3)^3$
\item $B = (1 + x)^3 - (2\sqrt{5} - 3x)(2\sqrt{5} + 3x)$
\item $C = (2x - 1)(x - 2)^3$
\item $D = (2x - 3)(4x^2 + 6x + 9)$
\item $E = (-2x + 3)^3 - 2x(2x^2 - 5x + 3)$
\item $F = (x - 2)^2(-2x^2 + x + 1)$
\end{enumerate}
\item Factoriser au mieux :
\begin{enumerate}
\item $F = 4x^2 + 20x + 25 - 9(4x^2 - 25)$
\item $G = 27 - 8x^3$
\item $H = 4(x - 3)^2 - (3x + 1)^2$
\item $I = x^3 - 8 + 2(x - 2)^2$
\end{enumerate}
\end{enumerate}
\end{multicols}

\fbox{\textbf{\exo}}
Exprimer sous forme de fractions irréductibles :
\begin{multicols}{2}
\setlength{\columnseprule}{0.1mm} % La largeur de la ligne verticale entre les colonnes
\begin{enumerate}
\item $A = \frac{(-3)^{2} \times 5^3 \times {10}^2}{2^3 \times 6 \times 3^6 \times (5 \times 2^{-3})^2}$
\item $B = \frac{(-6)^2 \times (7^{-2})^3 \times (9)^4}{(3^5) \times (-5)^{4} \times (-5)^{2}}$
\item $C = \frac{(-12)^3 \times (x^{-6})^{-3}}{(-9)^3 \times (x^{-3})^5}$
\item $D = \left( \frac{-27}{15}\right) ^{2} \times \left( \frac{(-3)^{2}}{5} \right)$
\item $E = \left( -3 \right)^4 \times \left( -5 \right) \times \left( -3 \right)^4 \times \frac{(-3)^{4}}{(-3)^{6}}\times\left[(-3)^{-2} \right]^{-1}$
\item $F = \frac{(-2)^5 \times  (-5^{8}) \times (-9^{3}) }{(-6)^4 \times (30)^5 }$
\item $G = (-2)^2 \times (-5) \times  \frac{(-3)^{3}}{5}$
\item $H = \left( \frac{-2}{3} +2 \right) \times \frac{3}{2}$
\end{enumerate}
\end{multicols}

\fbox{\textbf{\exo}}
simplifier l'écriture des réels suivants :
\begin{enumerate}
\begin{multicols}{2}
\item $A = 7\sqrt{20} - 11\sqrt{45} + 3\sqrt{80}$
\item $B = 2\sqrt{18} - 5\sqrt{8} + 4\sqrt{50}$
\item $C = 2\sqrt{27} - \sqrt{75} + 8\sqrt{12}$
\item $D = 2\sqrt{20} - 5\sqrt{45} + 6\sqrt{5} - 2\sqrt{180}$
\item $E = \frac{-3\sqrt{5}}{2\sqrt{15}}$
\item $F = \frac{2}{\sqrt{3} - 1}$
\item $G = \frac{\sqrt{3} - \sqrt{5}}{\sqrt{7} - 1} \times \frac{\sqrt{3} + \sqrt{5}}{\sqrt{7} + 1}$
\item $H = \frac{(\sqrt{3} - 2)^{2}}{\sqrt{7}} \times \frac{7 + 4 \sqrt{3}}{\sqrt{14}}$
\item $I = \frac{2}{1 - \sqrt{3}} + \frac{1}{1 + \sqrt{3}}$
\item $J = \frac{\sqrt{16}}{28} - \frac{\sqrt{125}}{49} - \frac{\sqrt{25}}{7}$
\end{multicols}
\item $K = \left( \sqrt{3} - 5 \right)^2 + (\sqrt{2} + 1)(\sqrt{2} - 1) + \sqrt{3}(5\sqrt{3} - 1)$
\item $L = (\sqrt{2} - 4)^2 + 5\sqrt{2}(3 + \sqrt{2}) - (7\sqrt{2} + 1)(7\sqrt{2} - 1)$
\end{enumerate}

\fbox{\textbf{\exo}}

\begin{enumerate}
\item Traduire les inégalités suivantes sous formes d'intervalles :
\begin{multicols}{3}
\setlength{\columnseprule}{0.1mm} % La largeur de la ligne verticale entre les colonnes
\begin{enumerate}
\item $2 < x \leq 4$
\item $-3 < x < 21$
\item $x < -5$
\item $x \geq -3$
\item $-4 < x < 25$
\item $x < 0$
\item $0 \leq x \leq 5$
\item $-1 < x \leq 15$
\item $x < 3$
\item $3x \geq 9$
\item $-2x < 8$
\end{enumerate}
\end{multicols}

\begin{multicols}{3}
\setlength{\columnseprule}{0.1mm} % La largeur de la ligne verticale entre les colonnes
\item Traduire les intervalles suivants sous formes d'inégalités :
\begin{enumerate}
\item $x \in ]-2; 5[$
\item $x \in ]-\infty; 5]$
\item $x \in [-3; 10[$
\item $x \in [-2; +\infty[$
\item $x \in ]11; 20[$
\item $x \in [0; +\infty[$
\item $x \in ]-\infty; 5]$
\end{enumerate}
\end{multicols}

\begin{multicols}{3}
\setlength{\columnseprule}{0.1mm} % La largeur de la ligne verticale entre les colonnes
\item Déterminer \( A \cap B \) et \( A \cup B \) dans chacun des cas suivants :
\begin{enumerate}
\item $A = [3; 10] \text{ et } B = [5; 7]$
\item $A = ]-\infty; 0] \text{ et } B = [0; +\infty[$
\item $A = [2; +\infty[ \text{ et } B = ]-5; 2]$
\item $A = [2; 11] \text{ et } B = ]-\infty; 0]$
\item $A = [4; 7] \text{ et } B = ]-\infty; 5]$
\item $A = ]-\infty; -3] \text{ et } B = [-3; +\infty[$
\end{enumerate}
\end{multicols}
\end{enumerate}

\fbox{\textbf{\exo}}
\begin{multicols}{3}
\setlength{\columnseprule}{0.1mm} % La largeur de la ligne verticale entre les colonnes
\begin{enumerate}
\item $|5 - x| - |-x + 7| = 0$
\item $|2x + 3| + 7 = 0$
\item $|8x - 1| = 5$
\item $|-x + 3| \leq 4$
\item $|-x + 6| + 8\geq 0$
\item $|2x + 5| > 7$
\item $|-3x - 2| < -9$
\item $|-2x + 1| - 5 \leq 0$
\item $|3x - 5| < 4$
\item $|-3x + 2| = |4x + 6|$
\item $|-2x + 3| - |3x - 5| = 0$
\end{enumerate}
\end{multicols}
\end{document}